% Appendix snapshot tables
\section{Repository Snapshot Tables}
\label{sec:appendix_snapshots}

This appendix records broader repository snapshot tables that are directionally consistent with the main text but were not part of the locked paper-evidence set when the quantitative claims in Section~\ref{sec:results} were written. We include them to document the current state of the project and to guide the next round of matched-leakage and deeper-layer analyses. The central paper claims continue to rely only on the completed paper suite summarized in the main text.

\subsection{Iterative Sycophancy Snapshot}
\begin{table}[t]
\centering
\caption{Repository snapshot from the iterative sycophancy report at $\lambda=1.0$ after 7 rounds, across layers 8, 16, and 23. Columns T (tangent-constrained) and A (ambient baseline) compare erasure methods. Utility is answer-matching delta relative to clean representations. Leakage and displacement decrease from left to right (lower is better); utility increases from left to right (higher is better). This snapshot table is not a main-text claim anchor and should be read only as evidence of directional consistency with the layer-8 paper-locked row (Table~\ref{tab:completed_results_summary}). Deeper-layer iterative tangent runs continue to reduce leakage more than ambient baselines across this exploratory sweep.}
\label{tab:repo_snapshot_sycophancy}
\small
\begin{tabular}{lcccccc}
\toprule
Layer & T leak & A leak & T util & A util & T disp & A disp \\
\midrule
8 & 0.063 & 0.213 & -20.7pp & -0.5pp & 2.88 & 0.62 \\
16 & 0.153 & 0.400 & +0.0pp & +0.0pp & 3.79 & 0.27 \\
23 & 0.138 & 0.361 & -0.1pp & -0.2pp & 44.50 & 3.11 \\
\bottomrule
\end{tabular}
\end{table}


\subsection{Iterative Safety Snapshot}
\begin{table}[t]
\centering
\caption{Repository snapshot from the iterative safety report at $\lambda=1.0$ after 7 rounds. Utility is pairwise helpfulness delta. These rows are directionally consistent with the main text but are not used as main-text claim anchors.}
\label{tab:repo_snapshot_safety}
\small
\begin{tabular}{lcccccc}
\toprule
Layer & T leak & A leak & T util & A util & T disp & A disp \\
\midrule
8 & 0.0725 & 0.2583 & -0.052 & -0.023 & 2.70 & 0.18 \\
16 & 0.1292 & 0.3517 & -0.092 & +0.003 & 4.37 & 0.60 \\
23 & 0.1483 & 0.3404 & -0.062 & -0.152 & 38.32 & 5.30 \\
\bottomrule
\end{tabular}
\end{table}


\subsection{Local CelebA Visual Table-2-Style Comparison}
\begin{table*}[t]
\centering
\caption{Local supplementary Table-2-style CelebA visual comparison using CLIP ViT-B/32 embeddings. For each target concept we intervene once per method and evaluate (i) concept leakage via retrained linear/nonlinear probes (lower is better) and (ii) mean control-task accuracy change across 3 least-correlated control attributes (closer to zero is better). Targets and controls: Male $\rightarrow$ Blurry,Narrow\_Eyes,Straight\_Hair; Smiling $\rightarrow$ Straight\_Hair,Bushy\_Eyebrows,Gray\_Hair; Blond\_Hair $\rightarrow$ Straight\_Hair,Narrow\_Eyes,Blurry. In this local supplementary setting, all methods use $\lambda{=}0.5$ except \TanGCE{} with $\lambda{=}8.0$.}
\label{tab:local_celeba_visual_table2}
\small
\setlength{\tabcolsep}{4pt}
\begin{tabular}{l|ccc|ccc|ccc}
\toprule
 & \multicolumn{3}{c|}{Male} & \multicolumn{3}{c|}{Smiling} & \multicolumn{3}{c}{Blond\_Hair} \\
Method & L$\downarrow$ & NL$\downarrow$ & $\Delta$Ctrl & L$\downarrow$ & NL$\downarrow$ & $\Delta$Ctrl & L$\downarrow$ & NL$\downarrow$ & $\Delta$Ctrl \\
\midrule
Random Global & .993 & .994 & -0.0pp & .891 & .911 & +0.0pp & .909 & .933 & -0.1pp \\
Random Per-Samp. & .993 & .994 & +0.0pp & .891 & .907 & +0.2pp & .907 & .932 & -0.1pp \\
Linear Proj. & .993 & .992 & +0.0pp & .900 & .906 & -0.3pp & .919 & .931 & -0.2pp \\
INLP & .994 & .993 & -0.0pp & .900 & .903 & +0.3pp & .918 & .932 & -0.0pp \\
LEACE & .994 & .991 & -0.0pp & .894 & .894 & +0.2pp & .909 & .925 & -0.1pp \\
\AmbGCE{} & .990 & .995 & -0.2pp & .858 & .884 & -0.4pp & .893 & .920 & +0.2pp \\
\TanGCE{} & .887 & .953 & -0.8pp & .649 & .649 & -0.6pp & .699 & .747 & -2.1pp \\
\bottomrule
\end{tabular}
\end{table*}


To ensure a modality-matched benchmark, this appendix now reports a vision-only CelebA run (CLIP ViT-B/32 embeddings) with the Table~\ref{tab:full_comparison} method family and 3 least-correlated controls per target concept. Artifacts are stored under \texttt{outputs/local/celeba\_visual\_table2\_v1/} (including per-method sample folders with 2 saved images each). The consolidated W\&B results run is \texttt{g52h8kbr}.

\subsection{CelebA Per-Concept Detail (All 40 Attributes)}
\label{sec:appendix_celeba_per_concept}
Tables~\ref{tab:celeba_clip_all40_per_concept_leak} and~\ref{tab:celeba_clip_all40_per_concept_ctrl} give the full per-concept breakdown behind the aggregate CelebA result in Table~\ref{tab:celeba_clip_all40}. The main-text Table~\ref{tab:celeba_clip_fair} reports three representative concepts (Male, Smiling, Blond\_Hair); the appendix tables cover all 40 binary CelebA attributes from the same \texttt{celeba\_clip\_all40\_fair\_v2} run under the shared validation-HPO protocol. TanGCE achieves the lowest nonlinear leakage on \textbf{36 of 40} concepts (LEACE wins on 3, IGBP on 1), with a mean margin of $0.116$ absolute NL accuracy over the best competing method per concept. Per-concept mean $\Delta$Ctrl remains within roughly a percentage point of zero, consistent with the aggregate mean NL of $0.707$ and mean $\Delta$Ctrl of $-0.7\%$ reported in the main text.
\begin{table*}[p]
\centering
\caption{Per-concept nonlinear probe leakage on CelebA (CLIP ViT-B/32), all 40 binary attributes. Each cell is retrained nonlinear probe test accuracy after erasure (lower = stronger erasure). $\lambda^*$ is selected per method per concept from $\{0.5, 1.0, 2.0, 4.0, 8.0\}$ by lowest validation NL accuracy. Bold marks the lowest leakage per row. Abbrev.: RandG = Random Global, RandS = Random Per-Sample, LinProj = Linear Projection, \AmbGCE{} = ambient gradient-based concept erasure, \TanGCE{} = tangent gradient-based concept erasure.}
\label{tab:celeba_clip_all40_per_concept_leak}
\footnotesize
\setlength{\tabcolsep}{3pt}
\begin{tabular}{l|cccccccc}
\toprule
Concept & RandG & RandS & LinProj & INLP & LEACE & IGBP & \AmbGCE{} & \TanGCE{} \\
\midrule
5\_o\_Clock\_Shadow & .92 & .92 & .91 & .92 & .87 & .92 & .91 & \textbf{.77} \\
Arched\_Eyebrows & .79 & .79 & .76 & .76 & .69 & .77 & .77 & \textbf{.59} \\
Attractive & .73 & .74 & .69 & .73 & .63 & \textbf{.53} & .73 & .60 \\
Bags\_Under\_Eyes & .76 & .75 & .74 & .75 & .70 & .77 & .75 & \textbf{.69} \\
Bald & 1.00 & .99 & .99 & .99 & .98 & 1.00 & .99 & \textbf{.97} \\
Bangs & .91 & .91 & .89 & .89 & .87 & .92 & .91 & \textbf{.68} \\
Big\_Lips & .80 & .79 & .74 & .80 & \textbf{.68} & .82 & .81 & .73 \\
Big\_Nose & .76 & .76 & .73 & .74 & .72 & .73 & .76 & \textbf{.61} \\
Black\_Hair & .85 & .84 & .82 & .81 & .79 & .83 & .82 & \textbf{.46} \\
Blond\_Hair & .92 & .92 & .94 & .92 & .92 & .92 & .92 & \textbf{.72} \\
Blurry & .98 & .98 & .98 & .96 & .97 & .98 & .97 & \textbf{.88} \\
Brown\_Hair & .80 & .80 & .76 & .75 & .75 & .80 & .80 & \textbf{.52} \\
Bushy\_Eyebrows & .87 & .88 & .86 & .85 & .76 & .87 & .86 & \textbf{.66} \\
Chubby & .94 & .94 & .95 & .95 & .92 & .95 & .94 & \textbf{.81} \\
Double\_Chin & .94 & .94 & .94 & .93 & .92 & .94 & .95 & \textbf{.85} \\
Eyeglasses & .99 & 1.00 & .96 & .93 & .96 & 1.00 & .96 & \textbf{.90} \\
Goatee & .96 & .97 & .97 & .93 & .96 & .96 & .95 & \textbf{.85} \\
Gray\_Hair & .97 & .97 & .97 & .97 & .97 & .97 & .98 & \textbf{.89} \\
Heavy\_Makeup & .89 & .89 & .89 & .86 & .85 & .88 & .88 & \textbf{.70} \\
High\_Cheekbones & .79 & .79 & .77 & .75 & .70 & .57 & .76 & \textbf{.50} \\
Male & .99 & .99 & .99 & .99 & 1.00 & .98 & 1.00 & \textbf{.85} \\
Mouth\_Slightly\_Open & .89 & .89 & .89 & .71 & .72 & .66 & .80 & \textbf{.43} \\
Mustache & .94 & .95 & .95 & .94 & .92 & .95 & .93 & \textbf{.77} \\
Narrow\_Eyes & .88 & .88 & .80 & .87 & .80 & .88 & .88 & \textbf{.78} \\
No\_Beard & .94 & .95 & .94 & .93 & .92 & .94 & .95 & \textbf{.77} \\
Oval\_Face & .70 & .71 & .69 & .68 & \textbf{.67} & .69 & .67 & .72 \\
Pale\_Skin & .93 & .93 & .94 & .95 & .93 & .92 & .93 & \textbf{.83} \\
Pointy\_Nose & .75 & .72 & .71 & .69 & \textbf{.63} & .72 & .75 & .64 \\
Receding\_Hairline & .94 & .94 & .96 & .96 & .94 & .94 & .95 & \textbf{.69} \\
Rosy\_Cheeks & .93 & .93 & .93 & .94 & .91 & .94 & .93 & \textbf{.79} \\
Sideburns & .95 & .95 & .94 & .94 & .95 & .94 & .94 & \textbf{.82} \\
Smiling & .90 & .89 & .91 & .74 & .82 & .59 & .88 & \textbf{.44} \\
Straight\_Hair & .78 & .78 & .79 & .77 & .78 & .78 & .78 & \textbf{.49} \\
Wavy\_Hair & .78 & .76 & .77 & .77 & .78 & .77 & .78 & \textbf{.52} \\
Wearing\_Earrings & .90 & .89 & .86 & .88 & .85 & .89 & .88 & \textbf{.61} \\
Wearing\_Hat & 1.00 & 1.00 & .97 & .97 & .97 & 1.00 & .97 & \textbf{.94} \\
Wearing\_Lipstick & .90 & .90 & .90 & .89 & .85 & .88 & .92 & \textbf{.75} \\
Wearing\_Necklace & .84 & .85 & .80 & .80 & .77 & .85 & .83 & \textbf{.68} \\
Wearing\_Necktie & .93 & .93 & .92 & .92 & .91 & .92 & .93 & \textbf{.78} \\
Young & .86 & .86 & .85 & .84 & .83 & .86 & .87 & \textbf{.60} \\
\midrule
\textbf{Mean} & .88 & .88 & .87 & .86 & .84 & .86 & .87 & \textbf{.71} \\
\bottomrule
\end{tabular}
\end{table*}

\begin{table*}[p]
\centering
\caption{Per-concept mean control accuracy change $\Delta$Ctrl (\%) on CelebA, all 40 binary attributes. For each concept, $\Delta$Ctrl is averaged over the 5 least-correlated control attributes (closer to zero = more surgical). $\lambda^*$ selection matches Table~\ref{tab:celeba_clip_all40_per_concept_leak}.}
\label{tab:celeba_clip_all40_per_concept_ctrl}
\footnotesize
\setlength{\tabcolsep}{3pt}
\begin{tabular}{l|cccccccc}
\toprule
Concept & RandG & RandS & LinProj & INLP & LEACE & IGBP & \AmbGCE{} & \TanGCE{} \\
\midrule
5\_o\_Clock\_Shadow & +0.0 & +0.1 & -0.7 & +0.1 & +0.0 & +0.1 & -2.3 & -0.2 \\
Arched\_Eyebrows & -1.3 & -2.3 & -1.2 & -0.8 & -1.1 & -2.5 & +0.6 & -0.6 \\
Attractive & -0.3 & -1.0 & -1.7 & -0.8 & -0.3 & -0.9 & +0.0 & -0.8 \\
Bags\_Under\_Eyes & -3.7 & -1.7 & -0.3 & -7.7 & -0.3 & -0.3 & -0.3 & -0.1 \\
Bald & +0.1 & +0.0 & -0.1 & +0.1 & -0.1 & +1.1 & +2.3 & +0.2 \\
Bangs & -3.7 & -1.2 & -0.5 & -7.9 & -0.1 & -0.7 & -0.3 & -0.8 \\
Big\_Lips & +0.1 & -0.0 & -0.3 & +0.0 & -0.1 & -0.5 & +0.1 & -0.1 \\
Big\_Nose & -0.1 & +0.3 & -3.5 & -9.7 & +0.0 & -0.3 & -1.8 & +0.1 \\
Black\_Hair & -0.1 & -0.0 & +0.1 & +0.2 & +0.3 & -0.3 & -0.3 & -0.7 \\
Blond\_Hair & -0.3 & +0.1 & -1.4 & +0.7 & -0.1 & -0.4 & -1.1 & -1.6 \\
Blurry & -1.3 & +0.0 & -0.1 & -5.5 & +0.1 & -0.7 & +0.7 & -0.1 \\
Brown\_Hair & +0.1 & -0.9 & -0.3 & -6.0 & -0.7 & +0.4 & +0.1 & -0.2 \\
Bushy\_Eyebrows & +0.0 & +0.1 & -0.1 & -1.5 & +0.2 & -0.9 & -0.3 & -0.5 \\
Chubby & -0.1 & +0.3 & -0.2 & -1.4 & -0.2 & -0.2 & +0.4 & -2.0 \\
Double\_Chin & +0.1 & +0.1 & -0.4 & -12.7 & +0.3 & +0.7 & -0.7 & -2.5 \\
Eyeglasses & -0.1 & +0.1 & -0.0 & -0.1 & -0.2 & -0.4 & -0.2 & +0.2 \\
Goatee & +0.1 & -0.1 & +0.0 & +0.2 & +0.1 & +0.5 & -0.3 & -0.4 \\
Gray\_Hair & -0.2 & -0.3 & -2.1 & +0.4 & -0.1 & -1.1 & -1.5 & -1.5 \\
Heavy\_Makeup & -0.3 & -0.1 & -1.1 & -1.7 & -1.3 & -1.3 & -1.1 & -0.9 \\
High\_Cheekbones & +0.1 & +0.2 & +0.1 & -0.3 & -0.1 & -0.4 & -0.0 & -0.7 \\
Male & -0.3 & -0.1 & -0.7 & -0.2 & -0.3 & -2.2 & -0.9 & -1.1 \\
Mouth\_Slightly\_Open & -0.3 & -0.6 & +0.1 & +0.4 & +0.2 & +0.3 & +0.1 & -0.2 \\
Mustache & +0.0 & -0.1 & +0.1 & -10.3 & -0.5 & +1.0 & +0.5 & -1.5 \\
Narrow\_Eyes & -0.1 & -0.2 & -0.3 & +0.2 & -0.1 & +0.3 & +0.0 & +0.0 \\
No\_Beard & -1.2 & -0.2 & -4.0 & -6.5 & -1.6 & -0.1 & -0.7 & -0.5 \\
Oval\_Face & -2.0 & +0.1 & -0.1 & -14.4 & -0.6 & -0.3 & +0.2 & -0.3 \\
Pale\_Skin & -1.2 & -0.1 & -0.1 & -6.8 & -0.9 & -0.1 & +0.7 & -1.0 \\
Pointy\_Nose & -0.1 & -0.2 & +0.3 & -1.1 & +0.2 & -0.7 & -0.1 & -0.1 \\
Receding\_Hairline & -0.1 & -0.1 & -0.1 & -5.9 & +0.0 & -0.3 & +0.1 & -1.1 \\
Rosy\_Cheeks & -0.5 & +0.1 & -0.6 & -0.1 & -0.2 & -0.1 & -0.1 & -1.3 \\
Sideburns & -0.4 & +0.0 & -0.1 & +0.1 & +0.1 & -0.6 & -2.5 & -0.1 \\
Smiling & -0.3 & -1.3 & +0.1 & +0.6 & +0.4 & +0.1 & +0.1 & -0.1 \\
Straight\_Hair & -0.1 & -0.2 & -0.1 & -0.9 & -0.1 & -0.8 & +0.5 & -0.7 \\
Wavy\_Hair & +0.1 & -0.7 & -0.9 & +0.2 & +0.1 & -0.4 & -0.1 & -0.7 \\
Wearing\_Earrings & +0.8 & -0.3 & -1.6 & -0.3 & -0.1 & -1.3 & -0.5 & -0.3 \\
Wearing\_Hat & -0.1 & +0.0 & -0.1 & -0.6 & +0.0 & +1.0 & +0.1 & -0.2 \\
Wearing\_Lipstick & +0.2 & -0.1 & -6.1 & -21.5 & -1.9 & -1.2 & -1.4 & -0.9 \\
Wearing\_Necklace & +0.1 & -1.7 & -1.7 & -9.3 & +0.3 & -1.8 & +0.2 & -1.2 \\
Wearing\_Necktie & +0.0 & +0.1 & +0.0 & +0.2 & -0.2 & -0.7 & -13.4 & -1.3 \\
Young & -0.0 & -0.7 & -1.7 & -10.7 & +0.1 & -0.9 & -2.2 & -0.8 \\
\midrule
\textbf{Mean} & -0.4 & -0.3 & -0.8 & -3.5 & -0.2 & -0.4 & -0.6 & -0.7 \\
\bottomrule
\end{tabular}
\end{table*}


\begin{table}[t]
\centering
\caption{Geometry evidence map organizing the appendix diagnostics into four research blocks. Each block addresses a distinct aspect of manifold structure: subspace sufficiency examines whether task-relevant computation concentrates in a low-dimensional linear subspace, perturbation sensitivity tests whether on-manifold directions matter more than orthogonal directions, representation connectivity evaluates whether the manifold remains connected across domains, and multi-layer accumulation tracks error compounding through the network. These diagnostics support the mechanism discussion and motivate the need for manifold-aware erasure methods but do not by themselves establish causality.}
\label{tab:geometry_schedule}
\scriptsize
\begin{tabular}{p{0.22\linewidth}p{0.22\linewidth}p{0.40\linewidth}}
\toprule
Experiment Block & Research Question & Primary Evidence \\
\midrule
Subspace sufficiency & Subspace concentration & Top-1 agreement under intrinsic projection \\
Perturbation sensitivity & High-variance vs. orthogonal directions & KL and top-1 under on/off-subspace noise \\
Representation connectivity & Shared vs. domain manifolds & Cross-domain interpolation smoothness \\
Multi-layer accumulation & Error compounding across the network & Simultaneous intrinsic projection across layers \\
\bottomrule
\end{tabular}
\end{table}


\paragraph{Geometry Note.}
The broader geometry report continues to support three useful points: early and mid-layer representations are effectively low-dimensional, on-manifold directions carry much of the model's sensitivity in these perturbation tests, and interpolation remains consistent with a connected representation space. At the same time, the Q2 perturbation study is best read as evidence of sensitivity concentration in a high-variance subspace rather than as a standalone proof of curved-manifold causality.
